\section{From arithmetic relations to Poisson laws}\label{sec:poisson-transfer}
The arithmetic estimate is now available independently of the probability
argument. The remaining difficulty is not a lack of pair information, but
the absence of a sparse unconditional dependency graph. Conditioning on the
small prime signs supplies exactly such a graph.

\subsection{An explicit transfer inequality}
Let $\mathcal F_Y=\sigma(X_p:p\le Y)$, with $Y>2B$. Let $D_Y$ be the
starts in $I_N$ for which some vertex of $V_x$ is $Y$-defective, and put
$G_Y=I_N\setminus D_Y$. We deliberately use the whole support in this
removal rule, so that the same convention works for marks and overlaps.
For $x\in G_Y$, the private-pivot lemma gives, for every environment,
\begin{equation}\label{eq:conditional-marginal}
 \PP(A_xX=b\mid\mathcal F_Y)=2^{-L}
 \quad\hbox{for every }b\in\F^L.
\end{equation}
After conditioning, an event at $x$ depends only on the prime coordinates
\[
 \Pi_x(Y)=\{p>Y:\exists n\in V_x,\ v_p(n)\hbox{ odd}\}.
\]
Join starts whose coordinate sets intersect. Disjoint collections with no
edge between them use disjoint independent coordinates. This is an exact
conditional dependency graph, not an approximate decorrelation assertion.

Write $E_Y(A)$ for its ordered edge count restricted to a mask $A$, and set
\[
 \mathcal M_B(A)=\sum_{x\in A}(2^{m_B(x)}-1),\qquad
 D_Y(A)=A\cap D_Y,\qquad \mathcal M_B=\mathcal M_B(I_N).
\]
All defect counts use the complete support, including $x-1$.
By \cref{cor:first-moment}, the actual probability mass removed from $A$
is at most $p(\mathcal M_B(A)+|D_Y(A)|)$, where $p=2^{-L}$.
Deleting the corresponding Poisson target coordinates costs at most
$p|D_Y(A)|$.
The edge count obeys
\begin{equation}\label{eq:edge-count}
 E_Y(I_N)\ll B^2\left(\frac{N^2}{Y\log Y}
                   +N\log\log N+\frac N{\log N}\right).
\end{equation}
Indeed a prime $p>Y$ can occur in at most $O(B(N/p+1))$ windows;
sum the square of this quantity over $Y<p\le3N$ and use partial summation.

We use the dependency-graph form of the Arratia--Goldstein--Gordon theorem
\cite{ArratiaGoldsteinGordon1989,ArratiaGoldsteinGordon1990}.
For indicators $(I_\alpha)$ with neighbourhoods $B_\alpha$ containing
$\alpha$, let
\[
 b_1=\sum_\alpha\sum_{\beta\in B_\alpha}p_\alpha p_\beta,
 \qquad
 b_2=\sum_\alpha\sum_{\substack{\beta\in B_\alpha\\\beta\ne\alpha}}
                  \EE(I_\alpha I_\beta),\qquad \mu=\sum_\alpha p_\alpha.
\]
An exact dependency graph has no third error term. The scalar
Poisson distance is at most $C(1\wedge\mu^{-1})(b_1+b_2)$; the joint
indicator-to-independent-Poisson distance is at most $C(b_1+b_2)$.
These assertions also apply on every fibre of $\mathcal F_Y$.

\begin{theorem}[Masked arithmetic-to-Poisson transfer]\label{thm:transfer}
For every deterministic $A\subset I_N$, put
$\mu_A=|A\cap G_Y|p$ and
$R_2(A)=\sum_{x,y\in A,\ |x-y|>L}(2^{\rho(x,y)}-1)$.
Then
\begin{align}
 &\EE\,\dTV\left(\Law\left(\sum_{x\in A}J_{x,L}\mid\mathcal F_Y\right),
                         \Pois(|A|p)\right)\notag\\
 &\quad\ll p(\mathcal M_B(A)+2|D_Y(A)|)
 +(1\wedge\mu_A^{-1})p^2\bigl(|A|+E_Y(A)+R_2(A)\bigr).
 \label{eq:transfer}
\end{align}
The multiplier is interpreted as one if $\mu_A=0$. Replacing it by one
gives a valid joint-field bound on the same finite site lattice.
\end{theorem}
\begin{proof}
Delete bad actual sites and their target coordinates. The natural coupling
costs at most $p(\mathcal M_B(A)+2|D_Y(A)|)$. On each remaining fibre, the
marginals are exactly $p$, and the graph above is a dependency graph.
Thus $b_1\le p^2(|A|+E_Y(A))$.
Overlapping starts are incompatible. At distance $L$, the union is a tree
with all its vertices possessing private coordinates above $Y$, since its
diameter is $2L<Y$. Its joint probability is therefore $p^2$.
For separated pairs the affine Fourier identity gives, after averaging,
\[
 \EE b_2\ll p^2\bigl(|A|+E_Y(A)+R_2(A)\bigr).
\]
Apply the scalar or process theorem on each fibre, average, and add the
deletions. The target is deterministic, so mixing also gives the same
unconditional bound.
\end{proof}
The denominator uses the mask's own mean. For a sparse mask it cannot be
replaced by the ambient intensity $Np$. At bounded ambient intensity,
all nonnegative sums are bounded by their full-block values, so the critical
estimate is uniform over all deterministic masks.

\subsection{The two prime-cutoff scales}\label{sec:cutoff-definition}
The leading scales are already visible without any special functions.
Set $S_N=\sqrt{\log N\log\log N}$ and $Y=e^{aS_N}$. The normalized number of bad
starts has leading upper exponent $1/(2a)$, while the normalized edge count
has exponent $a$. Their balance is $a=1/\sqrt2$.
For rare conditioning events and growing intensities we retain the smaller
second-order terms. The following definitions keep these terms explicit.

For $\nu\ge e$, let $u(\nu)\ge1$ solve $e^{u}/u=\nu$, and define
\[
 D(\nu)=\nu u(\nu)-\operatorname{Ei}(u(\nu)),\qquad
 H_N=\log N.
\]
For sufficiently large $N$, define $V_N$ and $\widehat V_N$ by
\begin{equation}\label{eq:saddle-definitions}
 V_N=D(H_N/V_N),\qquad
 \widehat V_N=2D(H_N/\widehat V_N),
 \qquad \nu_N=H_N/V_N,\quad\widehat\nu_N=H_N/\widehat V_N.
\end{equation}
The relevant solutions have $H_N/V_N,H_N/\widehat V_N\ge e$.
Since $D'(\nu)=u(\nu)>0$, each equation has exactly one such solution
for large $N$. Put $Y_N^\sharp=\lfloor e^{V_N}\rfloor$ and
$\widehat Y_N=\lfloor e^{\widehat V_N}\rfloor$.

\begin{proposition}[Uniform cutoff estimates]\label{prop:cutoffs}
Let $w\asymp\sqrt{H_N\log H_N}$, $\nu_w=H_N/w$ and $Y=\lfloor e^w\rfloor$.
Uniformly in a fixed logarithmic band,
\begin{equation}\label{eq:free-cutoff}
 \frac{|D_Y|}{N}\le e^{-D(\nu_w)+o(\nu_w)},\qquad
 \frac{E_Y(I_N)}{N^2}\le e^{-w+o(\nu_w)}+N^{-1+o(1)}.
\end{equation}
The maximum degree $\Delta_Y$ of the graph on all starts, including bad
ones, satisfies the same second bound after normalization by $N$.
Moreover,
\begin{align}
 V_N^2&=\frac{H_N}{2}
       \bigl(\log H_N+\log\log H_N-\log2-2+o(1)\bigr),
       \label{eq:hard-expansion}\\
 \widehat V_N^2&=H_N
       \bigl(\log H_N+\log\log H_N-\log4-2+o(1)\bigr).
       \label{eq:soft-expansion}
\end{align}
In particular $\widehat V_N\sim\sqrt2 V_N$ and
$\log\log N=o(\nu_N)$.
\end{proposition}
\begin{proof}
If $A(X,Y)$ counts $Y$-defective integers up to $X$, their unique
square--smooth factorization gives, for $0<\zeta<1/2$,
\[
 A(X,Y)/X\le X^{-\zeta}\prod_{p\le Y}(1+p^{-1+\zeta}).
\]
Choose $\zeta=u(\nu_w)/w$. The prime number theorem and partial summation
give $\sum_{p\le e^w}\log(1+p^{-1+\zeta})
=\operatorname{Ei}(u(\nu_w))+o(\nu_w)$.
This proves the first bound, because $|D_Y|\le B A(3N,Y)$ and
$\log B=o(\nu_w)$. The edge bound follows from \eqref{eq:edge-count}.
Each window uses $O(BH_N/w)$ primes above $Y$, each meeting
$O(B(N/Y+1))$ other windows, which proves the maximum-degree bound.
The expansions follow by substituting the asymptotic expansion of
$\operatorname{Ei}(u)$ into \eqref{eq:saddle-definitions}.
\Cref{supp:prop:cutoff-calculation} gives the uniform remainder calculation
and both expansions in detail.
\end{proof}

\begin{theorem}[Hard deletion and soft retention]\label{thm:scalar-rates}
Fix $0<\beta_-<\beta_+<\infty$. Uniformly when
$\beta_-\log N\le L+1\le\beta_+\log N$, for every $\varepsilon>0$,
\begin{align}
 \dTV(\Law(Z_{N,L}),\Pois(\lambda_N))
 &\ll\min\left\{1,\lambda_N
       \left(e^{-V_N+o(\nu_N)}+N^{-1/3+\varepsilon}\right)\right\},
       \label{eq:hard-rate}\\
 \dTV(\Law(Z_{N,L}),\Pois(\lambda_N))
 &\ll\min\left\{1,
       e^{-(\widehat V_N-\log^+\lambda_N)/2+o(\widehat\nu_N)}
                   +N^{-1/3+\varepsilon}\right\}.
       \label{eq:soft-rate}
\end{align}
The first estimate also bounds the mean conditional distance given
$\mathcal F_{Y_N^\sharp}$, and the second given
$\mathcal F_{\widehat Y_N}$. Neither saddle is asserted to be optimal for
the distribution itself.
\end{theorem}
\begin{proof}
At the hard cutoff, $|D_Y|=o(N)$ and $\mu_{I_N}\asymp\lambda_N$.
The master estimate and $Q_B=2N/\lambda_N$ give
\[
 R_2(N,L)/N^2\ll N^{-1/3+\varepsilon}(1+\lambda_N^{-1}).
\]
Its contribution to \eqref{eq:transfer} is
$O(\lambda_NN^{-1/3+\varepsilon})$. The defect mass contributes
$\lambda_NN^{-1/2+o(1)}$ and the cutoff terms contribute
$\lambda_Ne^{-V_N+o(\nu_N)}$. This proves \eqref{eq:hard-rate}.

For the second estimate retain the exceptional indicators in the scalar
sum instead of coupling all of them to zero. The scalar Stein equation
has factors $\|f\|_\infty\ll1\wedge\lambda^{-1/2}$ and
$\|\Delta f\|_\infty\le1\wedge\lambda^{-1}$
\cite{Krokowski2017,Rollin2012SteinFactors}.
Telescope only the good-start terms, using their exact conditional
marginal. Bound the exceptional part by its total first moment times
$\|f\|_\infty$. Exceptional neighbours of a good start stay in the
neighbour sums; their mean is controlled after averaging, not pointwise.
Set $\ell=\log^+\lambda$ and $K=\max(1,\lambda)=e^\ell$.
The same two Stein factors, at the same cutoff, give for every $\lambda>0$
the upper bound
\begin{equation}\label{eq:soft-free-rate}
 e^{\ell/2-D(H_N/w)+o(H_N/w)}
 +e^{\ell-w+o(H_N/w)}+K N^{-1/3+\varepsilon/2}.
\end{equation}
Here a smaller arithmetic exponent absorbs the subpolynomial terms.
The detailed calculation is in \cref{supp:sec:soft-application}.
At $w=\widehat V_N$ choose a common nonnegative remainder
$\eta_N\widehat\nu_N$, $\eta_N\to0$, in the two exponentials and put
$q=\exp(-(w-\ell)/2+\eta_N\widehat\nu_N)$.
If $q+N^{-1/3+\varepsilon}\ge1$, the bound $\dTV\le1$ suffices.
Otherwise $q<1$ implies $\ell\le w$, the second exponential is at most
$q$, and $K\le e^w\le N^{\varepsilon/2}$ for all sufficiently large $N$.
This proves \eqref{eq:soft-rate}. In particular intensities below one
are covered directly, without switching to a different conditioning
argument. All thresholds precede the length in the fixed band; the
fibrewise proof gives the claimed mean distance for the full
$\mathcal F_{\widehat Y_N}$.
\end{proof}
The hard estimate is more accurate at bounded intensity. The soft estimate
permits larger intensities: $\log^+\lambda_N\le\widehat V_N-c\widehat\nu_N$
for fixed $c>0$ suffices. In particular the hard-cutoff endpoint
$\log\lambda_N=V_N+O(1)$ is covered, but no convergence at the exact
new endpoint $\log\lambda_N=\widehat V_N$ is inferred.

Write $(z)_2=z(z-1)$ for the second falling factorial.
\begin{corollary}[First two moments]\label{cor:moments}
At critical lengths,
\begin{align*}
 \EE Z_{N,L}&=\lambda_N+O(N^{-1/2+o(1)}),\\
 \EE(Z_{N,L})_2&=\lambda_N^2+O_\varepsilon(N^{-1/3+\varepsilon}),\\
 \operatorname{Var}Z_{N,L}&=\lambda_N+O_\varepsilon(N^{-1/3+\varepsilon}).
\end{align*}
\end{corollary}
\begin{proof}
The first assertion is \cref{cor:first-moment}. Expand the second factorial
moment over ordered distinct starts. Overlaps cost $O(NLp^2)$;
touching pairs cost $p^2N^{1+o(1)}$ by \cref{lem:local-pairs}.
Separated pairs differ from $p^2$ by at most
$p^2(2^{\rho(x,y)}-1)$, whose sum is
$O_\varepsilon(N^{-1/3+\varepsilon})$. The variance identity concludes.
This argument proves the moments directly: total variation alone would
not imply convergence of unbounded moments.
\end{proof}
